\chapter{Introducci\'on}
\label{cap:introduccion}

La \escuela{} de la \universidad{} requiere definir un sistema que permita organizar el pr\'estamo de bienes de uso acad\'emico y recreativo bajo su custodia. El proyecto comprende exclusivamente libros, mesa de ping-pong, visor 3D, parlantes y carritos para rob\'otica que la Escuela autorice para este servicio. No comprende el pr\'estamo de computadoras ni la administraci\'on o el pr\'estamo de recursos de laboratorios.

El problema que se busca atender es la necesidad de conocer qu\'e bienes est\'an disponibles, qui\'en es responsable de cada pr\'estamo, cu\'ando debe devolverlos y qu\'e acciones corresponden ante un incumplimiento. Se plantea un registro centralizado de usuarios, inventario, reservas, entregas, devoluciones, garant\'ias, sanciones e historial. No se presupone un diagn\'ostico del proceso actual de la Escuela; este deber\'a contrastarse con sus responsables.

\section{Objetivo general}

Especificar los requisitos y el modelo de dominio de un sistema de gesti\'on de pr\'estamos de bienes para la Escuela, con control de disponibilidad, tiempos de uso, responsabilidades y trazabilidad.

\section{Objetivos espec\'ificos}

\begin{itemize}
    \item Identificar los actores y sus permisos: estudiante, docente y administrador.
    \item Delimitar el inventario prestable y diferenciar cada tipo de bien de sus unidades f\'isicas.
    \item Definir el ciclo de reserva, pr\'estamo, renovaci\'on y devoluci\'on.
    \item Especificar pol\'iticas configurables, garant\'ias y sanciones sin atribuir a la Escuela reglas a\'un no aprobadas.
    \item Establecer entidades, relaciones y restricciones que orienten el dise\~no posterior de la base de datos.
\end{itemize}

\section{Base del an\'alisis}

El an\'alisis de requisitos y modelo de dominio, e identifican bienes, usuarios, perfiles, tiempo de pr\'estamo, reservas, registros, inventario, estados del pr\'estamo, historial, garant\'ias, sanciones y pol\'iticas. 

Los detalles operativos que no aparecen en esos insumos se presentan como propuestas de dise\~no o decisiones pendientes de validaci\'on, no como normativa vigente. El informe se organiza en contexto, requisitos funcionales, requisitos no funcionales, modelo de dominio y conclusiones.
